\documentclass[12pt]{article}
\usepackage[utf8]{inputenc}
\usepackage[T1]{fontenc}
\usepackage[margin=1in]{geometry}
\usepackage{setspace}
\usepackage{xcolor}
\usepackage[colorlinks=true, urlcolor=blue!60!black, linkcolor=blue!60!black]{hyperref}
\setlength{\parindent}{2em}
\setlength{\parskip}{1.25em}
\renewcommand{\baselinestretch}{1.0}
\begin{document}
\begin{center}
{\LARGE \textit{United States v. Virginia}, 518 U.S. 515 (1996)\par}
\vspace{0.8em}
{\large Case Brief --- Four-Page Edition\par}
\end{center}
\vspace{1em}

\section{Memory Jogger}

The Court held 7--1 that Virginia violated the Equal Protection Clause by excluding women from the Virginia Military Institute without an ``exceedingly persuasive justification,'' and that a separate, inferior women's leadership program did not cure the violation (\href{https://supreme.justia.com/cases/federal/us/518/515/}{518 U.S. 515}).

\section{Detailed Case Facts}

VMI, founded in 1839 and supported and controlled by the Commonwealth, was the only single-sex school among Virginia's fifteen public colleges. Its 1,300 male cadets pursued a mission of producing ``citizen-soldiers'' through an ``adversative method'': physical rigor, mental stress, no privacy in spartan barracks, a seven-month ``rat line,'' and a strict honor code. VMI had the largest per-student endowment of any public undergraduate institution. In two years it received inquiries from 347 women and answered none. Prompted by a female high-school student's complaint, the United States sued in 1990.

After an appellate ruling against VMI, Virginia created the Virginia Women's Institute for Leadership (VWIL) at Mary Baldwin College, a private women's college. Designed around what was appropriate for ``most women,'' VWIL used a ``cooperative method which reinforces self-esteem,'' with no barracks and a ``largely ceremonial'' corps of cadets. Mary Baldwin's students averaged about 100 fewer SAT points, its faculty held fewer Ph.D.s, it offered no engineering degree, and its \$19 million endowment paled beside VMI's.

\section{Procedural History}

The district court ruled for Virginia, finding that single-sex education added diversity. The Fourth Circuit vacated, giving Virginia three options: admit women, create a parallel program, or privatize VMI. On remand the district court approved VWIL, and a divided Fourth Circuit affirmed under a new ``substantive comparability'' test, with Senior Judge Phillips calling VWIL a ``pale shadow'' of VMI. The Supreme Court granted certiorari and decided the case on June 26, 1996.

\section{Judicial Votes}

Liability affirmed and remedy reversed, \textbf{7--1}.
\begin{itemize}
\item \textbf{Majority:} Justice Ginsburg, joined by Justices Stevens, O'Connor, Kennedy, Souter, and Breyer.
\item \textbf{Concurring in the judgment:} Chief Justice Rehnquist.
\item \textbf{Dissenting:} Justice Scalia.
\item \textbf{Not participating:} Justice Thomas.
\end{itemize}

\section{Holding}

A State defending a sex classification must demonstrate an ``exceedingly persuasive justification,'' showing at least that the classification serves important governmental objectives and that the means are substantially related to them. The justification must be genuine, not \textit{post hoc}, and may not rest on overbroad generalizations about men and women. A separate program unequal in tangible resources and intangible prestige does not cure the exclusion of women from a unique public opportunity. Footnotes reserving strict scrutiny and limiting the decision to a ``unique'' opportunity are influential dicta.

\section{Analysis of Opinions}

\textbf{Majority (Ginsburg).} Tracing the Nation's ``long and unfortunate history of sex discrimination'' from \textit{Goesaert} to \textit{Reed}, Ginsburg described ``skeptical scrutiny.'' Sex classifications may compensate women for disabilities but never ``create or perpetuate the legal, social, and economic inferiority of women.'' Virginia's diversity rationale failed: VMI was founded when higher education was thought harmful to women, its other single-sex schools had gone coeducational, and a benefit for sons alone is not diversity. Fears that women would destroy the adversative method were ``self-fulfilling prophec[ies]''; some women could meet every demand, and the service academies had integrated. Under \textit{Sweatt v.\ Painter}, VWIL was a ``pale shadow'' of VMI in curriculum, faculty, funding, facilities, and prestige, and the Fourth Circuit erred in applying deferential review.

\textbf{Rehnquist (concurring in the judgment).} He objected to elevating ``exceedingly persuasive justification'' into the test, would not hold pre-1982 history against Virginia, and defined the violation as lacking an equal women's school, which Virginia could have cured without coeducating VMI.

\textbf{Scalia (dissenting).} Open, widespread traditions like all-male military colleges should be changed democratically. The Court effected a ``de facto abandonment'' of intermediate scrutiny, playing ``Supreme Court peek-a-boo'' with strict scrutiny; women, a majority of voters, are no ``discrete and insular minority.'' The Court ignored findings on single-sex pedagogy and left public single-sex education ``functionally dead.''

\section{Examples: Future Applications}

\begin{enumerate}
\item \textbf{Same-Side: Male-Only Maritime Academy.} A state academy excludes women, citing harsh shipboard discipline most women would not prefer. Generalizations about average preferences cannot bar qualified women.
\item \textbf{Same-Side: Underfunded Girls' STEM Wing.} After an all-boys magnet school is invalidated, a district opens a weaker girls' wing with fewer AP courses. This ``pale shadow'' remedy fails.
\item \textbf{Opposite-Side: Engineering Fellowship for Women.} A university funds fellowships for women in mechanical engineering, backed by data on underrepresentation. Compensatory classifications can survive.
\item \textbf{Opposite-Side: Police Academy Accommodations.} A coeducational academy separates showers and adjusts strength benchmarks. Footnote 19 approves privacy and physical-training adjustments.
\item \textbf{Fence-Sitter: Twin Single-Sex Schools.} A district simultaneously opens equal boys' and girls' high schools with identical funding and coeducational options. Footnote 7 reserved evenhanded single-sex offerings.
\end{enumerate}

\section{Critique}

\textbf{Originalist critiques.} Following Scalia, conservatives argue that the 1868 meaning did not forbid single-sex colleges and that \textit{Virginia} turned intermediate scrutiny into ``strict scrutiny in disguise.'' Scholarship on \textit{United States v.\ Skrmetti} (2025) calls Scalia's dissent the ``locus classicus'' of the critique, though Justice Barrett's \textit{Skrmetti} concurrence reaffirmed intermediate scrutiny.

\textbf{Progressive and feminist views.} Cass Sunstein called Ginsburg's opinion ``ambitious'' and ``nonminimalist,'' barring the conversion of difference into hierarchy. Mary Anne Case argued its demand for near-perfect proxies tracked post-\textit{Reed} precedent. Katharine Bartlett faulted the opinion for accepting VMI's hyper-masculine method as a baseline.

\section{Key Quotations}

\begin{quote}
``The justification must be genuine, not hypothesized or invented post hoc in response to litigation. And it must not rely on overbroad generalizations about the different talents, capacities, or preferences of males and females.''
\end{quote}
\noindent --- Justice Ginsburg, Opinion of the Court, 518 U.S. at 533.

\begin{quote}
``Parties who seek to defend gender-based government action must demonstrate an `exceedingly persuasive justification' for that action.''
\end{quote}
\noindent --- Justice Ginsburg, Opinion of the Court, 518 U.S. at 531.

\begin{quote}
``\,``Inherent differences'' between men and women, we have come to appreciate, remain cause for celebration, but not for denigration of the members of either sex or for artificial constraints on an individual's opportunity.''
\end{quote}
\noindent --- Justice Ginsburg, Opinion of the Court, 518 U.S. at 533.

\end{document}
